\documentclass[11pt]{article}
\usepackage[margin=1in]{geometry}
\usepackage[T1]{fontenc}
\usepackage[utf8]{inputenc}
\usepackage{lmodern,microtype,amsmath,amssymb,amsthm,mathtools}
\usepackage[colorlinks=true,linkcolor=blue,citecolor=blue,urlcolor=blue]{hyperref}
\hypersetup{pdftitle={Global-margin conditioning at the Gaussian endpoint},pdfauthor={Dmitry Vaintrob; Kaarel Hanni; OpenAI Codex}}
\setlength{\emergencystretch}{2em}
\newtheorem{theorem}{Theorem}[section]
\newtheorem{lemma}[theorem]{Lemma}
\newtheorem{proposition}[theorem]{Proposition}
\newtheorem{corollary}[theorem]{Corollary}
\theoremstyle{remark}\newtheorem{remark}[theorem]{Remark}
\DeclareMathOperator{\Var}{Var}
\DeclareMathOperator{\dist}{dist}
\DeclareMathOperator{\KL}{KL}
\DeclareMathOperator{\sign}{sign}
\newcommand{\E}{\mathbb E}
\newcommand{\Pp}{\mathbb P}
\title{Global-margin conditioning at the Gaussian endpoint:\\normalization, training loss, and output amplitude}
\author{Dmitry Vaintrob \and Kaarel H\"anni \and OpenAI Codex}
\date{October 15, 2026}
\begin{document}\maketitle
\begin{abstract}
We study the exact Gaussian-process posterior obtained by imposing a fixed absolute output margin at every Boolean input and then observing noiseless training signs. This is the $\gamma=0$ width limit of a specified Gaussian $\tanh$ network prior. For fixed depth, a random sample of size $n\le e^{d/25}$ has, simultaneously for every labeling, a raw conditioning probability at least $\exp[-C(d^2+nd+D_r\log d)]$, where $r\sim d/200$ and $D_r$ is the dimension of degree-$r$ Boolean polynomials. Consequently the actual latent posterior has average training log loss $O(d)$ when $n\ge D_r\log d$. We also prove that a posterior output draw has magnitude at least $c d^{(1-\beta)/2}$ at all but an $O(1/d)$ fraction of inputs in expectation, where $\beta=2^{-(L-1)}$. The proofs retain every global margin constraint. They use a comparison covariance, localized interpolation, and independent auxiliary subsets. These are conditioning estimates, not an exponential sample-complexity lower bound or an identification of posterior confidence with mean prediction. All kernel estimates needed for the conclusions are included.
\end{abstract}

\section{Model and results}
Let $\mathcal X=\{-1,1\}^d$, equipped with uniform probability, and $M=2^d$. Fix a depth $L\ge2$ and an absolute margin $a>0$. Define
\[
 \Psi(t)=\frac2\pi\arcsin t,\qquad F=\Psi^{\circ(L-1)},\qquad
 K(x,y)=F\!\left(\frac{1+x\cdot y}{d+1}\right),\qquad
 \beta=2^{-(L-1)}.
\]
For $G\sim N(0,K)$, inputs $S=(X_1,\ldots,X_n)$, and labels $y_i\in\{-1,1\}$, put
\[
 A_S=\{|G(x)|>a\ \forall x\in\mathcal X;\ y_iG(X_i)>0\ \forall i\},
 \qquad D(S)=-\log\Pp(A_S).
\]
The posterior throughout this paper is the law of $G$ conditional on $A_S$. This is exactly the same law as first conditioning on the global margin and then on training signs. The probability in $D(S)$ is taken under the original, unconditioned Gaussian law.

The Walsh character $\chi_U(x)=\prod_{i\in U}x_i$ is an eigenfunction of the normalized covariance operator $K/M$. Its eigenvalue is denoted by $\lambda_{|U|}$. For an odd cutoff $r$ let $D_r=\sum_{j\le r}\binom dj$ and let $\tau_r$ be the diagonal of the covariance retaining degrees greater than $r$. Write $v=\lambda_{\min}(K)>0$. The distinction between the matrix eigenvalue $v$ and the normalized eigenvalues $\lambda_j$ is essential.

There is an exact splitting
\[
 G=H+\sqrt v\,\xi,\qquad H\sim N(0,K-vI_M),\quad \xi\sim N(0,I_M),
\]
with independence. The latent covariance may be singular. Write $\rho_S$ for the actual conditional law of $H$ given $A_S$, and define
\[
 B_a(t)=\Phi\!\left(\frac{t-a}{\sqrt v}\right)+\Phi\!\left(\frac{-t-a}{\sqrt v}\right),\quad
 p_a(t)=\frac{\Phi((t-a)/\sqrt v)}{B_a(t)},\quad \ell_a(t)=-\log p_a(t).
\]
The quantity $p_a(t)$ is the probability of a positive output conditional on latent value $t$ and on its absolute margin.

\begin{theorem}\label{thm:main}
Fix $L\ge2$ and $a>0$, and let $r$ be the largest odd integer not exceeding $d/200$. There are positive constants $C,c,b$, depending only on $L,a$, such that the following holds for sufficiently large $d$. For every integer $1\le n\le e^{d/25}$, an independent uniform input sample has probability at least $1-e^{-cd}$ of satisfying, simultaneously for every labeling,
\begin{align}
 D(S)&\le C(d^2+nd+D_r\log d),\label{eq:main-denominator}\\
 \E_{\rho_S}\frac1n\sum_{i=1}^n\ell_a(y_iH(X_i))
 &\le C\left(d+\frac{d^2+D_r\log d}{n}\right),\label{eq:main-loss}\\
 \E_{\rho_S}\frac1n\sum_{i=1}^n(y_iH(X_i))_-^2
 &\le Cv\left(d+\frac{d^2+D_r\log d}{n}\right),\label{eq:main-wrong}\\
 \E\!\left[\Pp_x\{|G(x)|\le b d^{(1-\beta)/2}\}\mid A_S\right]
 &\le \frac Cd.\label{eq:main-amplitude}
\end{align}
Here $t_- =\max\{-t,0\}$. Each specified sample size has the stated probability guarantee; no simultaneous assertion over all prefixes of one infinite sample is intended.
\end{theorem}

Since $D_r\log d$ eventually exceeds $d^2$, \eqref{eq:main-loss} is $O(d)$ when $n\ge D_r\log d$. Estimate \eqref{eq:main-amplitude} concerns the actual output draw, averaged over the posterior and a uniform input. It does not assert a bound at every input for every posterior draw.

\subsection{Relation to the network convention}
For completeness, consider width-$N$ networks
\[
 h^0(x)=(1,x),\quad h^s(x)=\tanh(W^sh^{s-1}(x))\ (1\le s<L),\quad z_N(x)=W^Lh^{L-1}(x),
\]
with independent centered Gaussian entries and variances
\[
 \Var W^1_{ij}=N/d,\qquad \Var W^s_{ij}=1\ (1<s<L),\qquad \Var W^L_i=N^{-1}.
\]
These are the $\gamma=0$ scalings. At fixed $d,L$, the empirical hidden-layer covariances converge, by conditional laws of large numbers and induction, to iterates of $\Psi$. Indeed preactivation variances diverge, $\tanh$ converges to the sign function under a Gaussian coupling, and the Gaussian sign-correlation identity is $\E\sign U\sign V=(2/\pi)\arcsin\rho$. Continuous matrix square roots and bounded convergence justify the conditional covariance limit, including singular intermediate matrices. The first layer normalizes the augmented inputs to correlation $(1+x\cdot y)/(d+1)$. The final output is conditionally Gaussian with the empirical last-hidden covariance, so $z_N\Rightarrow G$ as a vector on the finite cube.

The spectral proof in Appendix~\ref{app:spectrum} shows $K$ is positive definite. Every boundary $G(x)=\pm a$ or $G(x)=0$ is a null set, and every compatible open margin-and-label event has positive probability. Ratios of probabilities therefore show convergence to the posterior used here when $N\to\infty$ at fixed $d,L,a,S$. No interchange of width and dimension limits is claimed.

\section{A comparison covariance and a global interpolation bound}
For a covariance $J$ on the cube, use the Gaussian Cameron--Martin norm $\|h\|_J^2=h^TJ^{-1}h$. For a translation-invariant kernel it equals $\sum_U\widehat h(U)^2/\lambda^J_{|U|}$, using normalized Walsh coefficients.

Let $Q$ retain degrees greater than $r$, let $K_{\rm low}=K-Q$, and set $\tau=\tau_r$. Define
\[
 \widetilde K=\tau K_{\rm low}+Q.
\]
Appendix~\ref{app:spectrum} proves $\tau\asymp d^{-\beta}$, $\lambda_0\ge c/d$, and $v\ge c(d\log d)^{-\beta}$. In particular $0<\tau<1$ for large $d$. The diagonal of $\widetilde K$ is $\tau(1-\tau)+\tau\le2\tau$. Its low-degree eigenvalues are those of $K$ multiplied by $\tau$, while all remaining eigenvalues agree. The Gaussian densities thus satisfy, pointwise,
\begin{equation}\label{eq:density-compare}
 \frac{dN(0,K)}{dN(0,\widetilde K)}(g)
 =\tau^{D_r/2}\exp\!\left[\frac12g^T(\widetilde K^{-1}-K^{-1})g\right]
 \ge\tau^{D_r/2}.
\end{equation}
The comparison covariance is only a proof device; it does not replace the prior in the theorem.

\begin{lemma}[Localized interpolation]\label{lem:bump}
With probability at least $1-e^{-cd}$ over the sample in Theorem~\ref{thm:main}, simultaneously for every labeling and every $A>0$, there is a function $h$ such that
\[
 |h(x)|>A\quad\forall x,\qquad y_i h(X_i)>A\quad\forall i,
 \qquad \|h\|_{\widetilde K}^2\le C A^2(d/\tau+nd^\beta).
\]
\end{lemma}
\begin{proof}
For $z\in\mathcal X$ define
\[
 b_z(x)=20\int_{2/5}^{9/20}(1-2p)^{\dist(x,z)}\,dp.
\]
Then $b_z(z)=1$, $0\le b_z(x)\le(1/5)^{\dist(x,z)}$, and its Walsh coefficients have magnitudes
\[
 \widehat b_j=20\int_{2/5}^{9/20}(1-p)^{d-j}p^j\,dp,
 \qquad \binom dj\widehat b_j\le\frac{20}{d+1},\qquad \widehat b_j\le(3/5)^d.
\]
The middle inequality follows by extending to the complete beta integral. Appendix~\ref{app:spectrum} gives
\[
 \lambda_j\ge\frac{c d^{-1-\beta}}{\binom dj}\quad(j\le .49d).
\]
Therefore the contribution of these degrees to $\|b_z\|_K^2$ is at most
\[
 C d^{1+\beta}\sum_{j\le .49d}\left[\binom dj\widehat b_j\right]^2\le C d^\beta.
\]
For $j>.49d$ the integrand increases with $p$ on the integration interval, giving $\widehat b_j\le .55^{d-j}.45^j$. Since every $\lambda_j\ge v/M$, the remaining energy is at most
\[
 \frac{2^d}{v}\sum_{j>.49d}\binom dj .55^{2(d-j)}.45^{2j}\le \operatorname{poly}(d)e^{-c_*d},
\]
where
\[
 c_*=-[\log2+h(.49)+1.02\log(.55)+.98\log(.45)]>.0062.
\]
Here $h(t)=-t\log t-(1-t)\log(1-t)$, and the binomial entropy bound applies because the exponent decreases for $j/d\ge .49$.

The low-degree energy increases by at most
\[
 \tau^{-1}D_r( .72)^d/v
\]
under damping. Since $D_r\le\exp[(h(1/200)+o(1))d]$ and $h(1/200)<.04$, this term decays exponentially. Hence $A_d:=\|b_z\|_{\widetilde K}^2\le Cd^\beta$.

Let $J(z,w)=\langle b_z,b_w\rangle_{\widetilde K}$. Its normalized Fourier eigenvalues are $\nu_j=\widehat b_j^2/\lambda^{\widetilde K}_j\le .72^d/(\tau v)$. Parseval and Cauchy--Schwarz give
\[
 \E_{X,Y}|J(X,Y)|\le\left(\sum_j\binom dj\nu_j^2\right)^{1/2}
 \le\sqrt{\nu_{\max} A_d}.
\]
For every subset $T$ of the sampled indices, even one chosen after seeing the sample,
\[
 \left\|\sum_{i\in T}b_{X_i}\right\|_{\widetilde K}^2
 \le nA_d+\sum_{i\ne j}|J(X_i,X_j)|.
\]
The expected cross sum is at most $n(n-1)\sqrt{\nu_{\max}A_d}$. Markov's inequality at threshold $nd^\beta$ has exponentially small failure when $n\le e^{.04d}$, because $\tfrac12\log(1/.72)>.16$. Thus all subsets simultaneously have energy at most $Cnd^\beta$.

Also, with exponentially high probability, every pair of different sample indices has inputs at distance greater than $.2d$: a union bound costs at most $n^2\exp[-d(\log2-h(.2))]$, with $\log2-h(.2)>.19$. This event excludes repetitions. At any input at most one center is within distance $.1d$, and the contribution of all other centers is at most $n(.2)^{.1d}=o(1)$. Thus the bump sum over any selected centers is $1+o(1)$ at a selected center and at most $.2+o(1)$ everywhere else.

For $T=\{i:y_i=-1\}$ take $h=2A(1-2\sum_{i\in T}b_{X_i})$. The preceding bounds give its strict global and training margins. Since $\|1\|_{\widetilde K}^2=1/(\tau\lambda_0)\le Cd/\tau$, the squared triangle inequality gives the asserted energy bound.
\end{proof}

\begin{proof}[Proof of \eqref{eq:main-denominator}]
Put $t_0=\sqrt{4\tau\log(4M)}$ and $A=a+t_0+1$. Under $N(0,\widetilde K)$ the symmetric event $B=\{\max_x|G(x)|\le t_0\}$ has probability at least $1/2$, by the diagonal bound and a Gaussian union bound. Lemma~\ref{lem:bump} makes $h+B$ a subset of $A_S$. Because $a$ is fixed, $\tau\asymp d^{-\beta}$, and $\beta<1$, we have $A^2\le C_a\tau d$ for large $d$. Hence
\[
 \|h\|_{\widetilde K}^2\le C(d^2+nd).
\]
For a symmetric event $B$, the Gaussian translation formula, symmetry, and $\cosh t\ge1$ imply
\[
 \Pp_{\widetilde K}(h+B)\ge \Pp_{\widetilde K}(B)e^{-\|h\|_{\widetilde K}^2/2}.
\]
Combine this with \eqref{eq:density-compare} and $\log(1/\tau)=O(\log d)$ to obtain \eqref{eq:main-denominator}.
\end{proof}

\section{The exact posterior training loss}
On the sample event in Lemma~\ref{lem:bump}, all observed inputs are distinct. Conditional independence of the noise coordinates in $G=H+\sqrt v\xi$ gives the exact latent posterior density
\[
 d\rho_S(H)=\frac{\exp[-V(H)-\sum_i\ell_a(y_iH(X_i))]}{\Pp(A_S)}\,d\nu(H),
 \qquad V(H)=-\sum_{x\in\mathcal X}\log B_a(H(x)),
\]
where $\nu=N(0,K-vI_M)$. All terms in the exponent are nonnegative. The Gibbs identity is consequently
\begin{equation}\label{eq:gibbs}
 D(S)=\KL(\rho_S\Vert\nu)+\E_{\rho_S}V(H)+\E_{\rho_S}\sum_i\ell_a(y_iH(X_i)).
\end{equation}
This holds on the support of $\nu$ even when its covariance is singular. The density is bounded above by $e^{D(S)}$, which also justifies integrability and the identity. Nonnegativity proves \eqref{eq:main-loss}.

For completeness, $\ell_a(t)\ge(t_-)^2/(2v)$. To see this write $t=-u\le0$ and $\overline\Phi=1-\Phi$. If $0\le u\le a$, the normal-tail ratio bound
\[
 \frac{\overline\Phi(A)}{\overline\Phi(B)}\le e^{-(A^2-B^2)/2}\qquad(A\ge B\ge0)
\]
implies $p_a(-u)\le e^{-2au/v}$. The ratio bound follows by integrating the normal hazard $\phi(t)/\overline\Phi(t)\ge t$. If $u\ge a$, the denominator in $p_a(-u)$ is at least $1/2$ and $\overline\Phi(t)\le\tfrac12e^{-t^2/2}$ for $t\ge0$, giving $p_a(-u)\le e^{-(a+u)^2/(2v)}$. Both cases imply the claimed inequality; for $t\ge0$ it is immediate. Equation~\eqref{eq:gibbs} now proves \eqref{eq:main-wrong}.

\section{Auxiliary subsets force large actual output magnitudes}
The next lemma is independent of the cube and of the training likelihood.

\begin{lemma}[A finite Gaussian amplitude estimate]\label{lem:auxiliary}
Let $G$ be a nondegenerate centered Gaussian vector on a finite set. Suppose an event $A$ implies $|G(x)|>a$ at every coordinate and $-\log\Pp(A)\le D$, with $D\ge1$. Let $T=(x_1,\ldots,x_m)$ be distinct coordinates with $K_T\succeq wI_m$, where $w>0$. Assume $D<m$ and $e^{-2a^2/w}\le1/2$. Set
\[
 t=\sqrt{\log(m/D)},\qquad b_0=a+\sqrt w\,t,\qquad c_0=\frac{e^{-3/2}}{\sqrt{2\pi}}.
\]
Then
\begin{align}\label{eq:auxiliary}
 &\E\!\left[\frac1m\sum_{i=1}^m\mathbf1_{\{|G(x_i)|\le b_0/2\}}\,\middle|\,A\right]\\
 &\qquad\le \frac2{c_0}(t+1)\frac{D+1}{\sqrt{mD}}+\frac{8w}{b_0^2}+\frac{16wD}{mb_0^2}.\notag
\end{align}
The event may impose arbitrary further constraints outside $T$.
\end{lemma}
\begin{proof}
On an enlarged Gaussian space draw $H_T\sim N(0,K_T-wI)$ and $\xi\sim N(0,I_m)$ independently, set $G_T=H_T+\sqrt w\xi$, and draw the other coordinates from their original conditional law given $G_T$. This reproduces the original full law of $G$. The outside coordinates may depend on $\xi$; no independence from them is asserted. The event inclusion alone gives
\[
 \Pp(A\mid H_T)\le\prod_{i=1}^m B_{a,w}(H_T(i)),\quad
 B_{a,w}(h)=\Phi((h-a)/\sqrt w)+\Phi((-h-a)/\sqrt w).
\]
The even function $B_{a,w}$ increases with $|h|$. If $q_H=m^{-1}\#\{i:|H_T(i)|\le b_0\}$ and $c_b=-\log B_{a,w}(b_0)$, then
\[
 \Pp(q_H\ge q\mid A)\le\min\{1,e^{D-mc_bq}\},\qquad
 \E[q_H\mid A]\le\frac{D+1}{mc_b}.
\]
Now $1-B_{a,w}(b_0)=\overline\Phi(t)-\overline\Phi(t+2a/\sqrt w)\ge\overline\Phi(t)/2$. The ratio bound used above proves this last inequality. Integrating the normal density over $[t,t+1/(t+1)]$ gives $\overline\Phi(t)\ge c_0 e^{-t^2/2}/(t+1)$. Since $-\log B\ge1-B$, these estimates give the first term on the right of \eqref{eq:auxiliary} as an upper bound for $\E[q_H\mid A]$.

The enlarged posterior has relative entropy exactly $-\log\Pp(A)\le D$ from its prior. Data processing gives the same upper bound for the relative entropy of its $\xi$ marginal from $N(0,I_m)$. The entropy variational inequality and the Gaussian moment generating function imply
\[
 \tfrac14\E[\|\xi\|^2\mid A]\le D+\log\E e^{\|\xi\|^2/4}=D+\tfrac m2\log2,
\]
so $\E[\|\xi\|^2\mid A]\le4D+2m$. If $|G_T(i)|\le b_0/2$ and $|H_T(i)|>b_0$, then $\sqrt w|\xi_i|>b_0/2$. The averaged Markov inequality completes the proof.
\end{proof}

\begin{proof}[Proof of \eqref{eq:main-amplitude}]
Hold fixed any training sample on the event proving \eqref{eq:main-denominator}, and any labeling. Increase its deterministic upper bound $D$ to be at least one. Since $n\le e^{.04d}$ and $h(1/200)<.04$, for large $d$ we can take $D\le e^{.045d}$.

Independently draw an auxiliary uniform sample $T$ of size $m=\lfloor e^{.05d}\rfloor$. Let $u$ be odd with $u/d\to1/50$, and put $w=\tau_u/2\asymp d^{-\beta}$. Appendix~\ref{app:gram} proves, with exponentially high probability over $T$,
\[
 Q_{>u,TT}\succeq (\tau_u/2)I_m,
\]
and thus $K_T\succeq wI_m$; this event in particular excludes repeated points. Indeed the bounds there decay exponentially because
\[
 .05<h(1/50),\qquad .10<\log2-h(1/4).
\]
The far-pair error is exponentially small as well, since $u$ is a fixed positive fraction of $d$.

Lemma~\ref{lem:auxiliary} applies for every such $T$, to the same full event $A_S$. Here $\log(m/D)\ge .004d$ for all sufficiently large $d$, so $b_0^2\ge c d^{1-\beta}$; its first error term is exponentially small and the other terms are $O(1/d)$. Choose a fixed positive $b$ with $b d^{(1-\beta)/2}\le b_0/2$.

Although the enlarged splitting in the lemma depends on $T$, its marginal law of the full $G$ is always the original Gaussian law, and its conditional marginal is always the same actual posterior $G\mid A_S$. Averaging its conclusion over the independent auxiliary sample therefore changes the left side exactly into
\[
 \E[\Pp_x\{|G(x)|\le b d^{(1-\beta)/2}\}\mid A_S].
\]
Bad auxiliary samples contribute at most their exponentially small probability. This proves \eqref{eq:main-amplitude} and completes Theorem~\ref{thm:main}.
\end{proof}

\section{What these estimates do not prove}
For parity labels $f(x)=\prod_i x_i$, define the actual confidence $q_S(x)=\Pp(f(x)G(x)>0\mid A_S)$ and the failure fraction $F_S=\Pp_x\{q_S(x)\le2/3\}$. Theorem~\ref{thm:main} does not give a positive lower bound on $F_S$. Its training-loss bound still needs a valid comparison with population loss for this globally conditioned posterior. Its amplitude estimate gives magnitudes, not signs or correlation with parity.

In particular, a small positive lower bound on the posterior-draw error would be insufficient: if $R_S=\E[\Pp_x(f(x)G(x)<0)\mid A_S]$, then $R_S\le1/3+(2/3)F_S$. This implication would require $R_S>1/3$ to force a positive confidence-failure fraction. Nor is the sign of a posterior mean used as a substitute for confidence. The latent posterior contains the full factor $\prod_xB_a(H(x))$; no log-concavity of this posterior is assumed. Establishing an exponential sample lower bound with these conventions remains outside the conclusions of this note.

\appendix
\section{Fourier level masses of the kernel}\label{app:spectrum}
Put $D=d+1$, $\beta=2^{-(L-1)}$, and $F=\Psi^{\circ(L-1)}$.
The characters $\chi_U(x)=\prod_{j\in U}x_j$ diagonalize the
normalized covariance operator $M^{-1}K$, $M=2^d$. Let $\lambda_j$
be its eigenvalue at degree $j$. We first work on the full $D$-cube
with kernel $F(u\cdot v/D)$, whose degree-$j$ eigenvalue is denoted
by $a_j$. Constants in this and the next section may depend on $L$.
Notation $b\asymp c$ means two inequalities by positive constant
multiples, uniformly in the stated ranges.

\begin{lemma}[Power coefficients and paired eigenvalues]\label{lem:power}
Write $F(z)=\sum_{m\text{ odd}}c_mz^m$. Every odd coefficient is
positive, $\sum_m c_m=1$, and
\[
 c_m\asymp\frac{\Gamma(m-\beta)}{m!}\asymp m^{-1-\beta}
 \qquad(m\ge1\text{ odd}).
\]
The full-cube eigenvalues vanish at even degrees and strictly decrease
through odd degrees. They are positive at every odd degree. On the
original $d$-cube,
\[
 \lambda_{2k}=\lambda_{2k+1}=a_{2k+1}>0,
\]
with a final unpaired eigenvalue if $d$ is even. In particular $K$
is positive definite.
\end{lemma}
\begin{proof}
The series of $\Psi$ has positive coefficients at every odd exponent,
with
\[
 [z^{2k+1}]\Psi(z)=\frac2\pi\frac{\binom{2k}{k}}{4^k(2k+1)}.
\]
Composition preserves strict positivity and oddness. Near $z=1$ in
a slit neighborhood,
\[
 \Psi(z)=1-b(1-z)^{1/2}+O((1-z)^{3/2}),
 \qquad b=\frac{2\sqrt2}{\pi}.
\]
Induction gives
$F(z)=1-A(1-z)^\beta+O((1-z)^{3\beta})$, with $A>0$;
at each composition its leading constant changes from $A_t$ to
$b\sqrt{A_t}$, starting from $A_1=b$.
These expansions admit algebraic coefficient transfer
\cite[Chapter VI]{flajolet2009}. Here are the analytic-domain details.
On the unit circle outside $\{-1,1\}$, positivity of every odd
coefficient gives $|\Psi(z)|<1$: equality in the triangle inequality
would require $z^2=1$. Hence the iterated functions extend
analytically through every compact arc away from these endpoints.
At $1$, restrict $\arg(1-z)$ to a closed proper subinterval of
$(-\pi,\pi)$. Each square-root composition reduces this angular
range, and its positive leading coefficient keeps its image in the
slit neighborhood for the next arcsine composition. Oddness gives
the identical argument at $-1$. These local sectors and the compact
arcs provide the required two-singularity indented domain.
The contributions of $1$ and $-1$ add at odd exponents, giving
\[
 c_m\sim\frac{2A}{-\Gamma(-\beta)}m^{-1-\beta}
 \quad(m\to\infty\text{ odd}).
\]
Together with positivity of the finitely many initial coefficients
and the gamma-ratio asymptotic, this proves the coefficient comparisons.
For $L=2$ the explicit formula and Stirling's bounds suffice instead.

Let $p_m(j)$ be the probability that the simple coordinate-flip walk
on the $D$-cube reaches a specified distance-$j$ vertex in $m$ steps.
Expanding $(u\cdot v/D)^m$ shows that its degree-$j$ Fourier coefficient
is $p_m(j)$. Thus $a_j=\sum_{m\text{ odd}}c_mp_m(j)$.
For the walk recursion,
\begin{align*}
 D[p_{m+1}(j)-p_{m+1}(j+2)]
 ={}&j[p_m(j-1)-p_m(j+1)]\\
 &+(D-j-2)[p_m(j+1)-p_m(j+3)].
\end{align*}
Induction from time zero gives nonincreasing probabilities through
each parity class; terms with zero boundary coefficients are omitted.
The term $m=j$ makes the odd-eigenvalue inequalities strict and proves
positivity. Even-degree eigenvalues vanish by parity.
Restricting the distinguished coordinate to $+1$ gives
$\lambda_q=a_q+a_{q+1}$, proving the claims.
\end{proof}

For odd $r$, the exact combined multiplicity of reduced degrees
$j-1,j$ is $\binom d{j-1}+\binom dj=\binom Dj$. Consequently
\begin{equation}\label{eq:paired}
 D_r=\sum_{\substack{j\le r\\j\text{ odd}}}\binom Dj,
 \quad
 \tau_r=\sum_{\substack{j>r\\j\text{ odd}}}\binom Dj a_j,
 \quad \lambda_{r+1}=a_{r+2}.
\end{equation}
The odd cutoff is essential for these exact identities.

\begin{lemma}[Sharp level and tail estimates]\label{lem:sharp}
Uniformly for odd $1\le j\le D/4$ and odd $1\le r\le D/8$,
\begin{align}
 \binom Dj a_j&\asymp j^{-1-\beta},\label{eq:level}\\
 \tau_r&\asymp r^{-\beta},\label{eq:trace}\\
 \frac{\lambda_{r+1}}{\tau_r}
 &\asymp\frac1{r\binom D{r+2}},\label{eq:ratio}\\
 s_{2,r}:=\sum_{q>r}\binom dq\lambda_q^2
 &\le\frac{C r^{-2\beta-2}}{\binom D{r+2}}.
 \label{eq:secondmoment}
\end{align}
The same level-mass estimate holds up to $j\le\rho D$ for any fixed
$\rho<1/2$, with constants depending on $\rho$.
\end{lemma}
\begin{proof}
The gamma integral and Lemma~\ref{lem:power} give, for odd $j$,
\begin{align}
 a_j&\asymp D^{-\beta}\int_0^\infty
 u^{-1-\beta}p(u)^j(1-p(u))^{D-j}\,du,
 \quad p(u)=\frac{1-e^{-2u}}2,\label{eq:poisson}\\
 &\asymp D^{-\beta}\int_0^{1/2}p^j(1-p)^{D-j}w_\beta(p)\,dp,
 \label{eq:betaweight}
\end{align}
where $w_\beta(p)=[-\frac12\log(1-2p)]^{-1-\beta}/(1-2p)$.
Indeed the gamma integral Poissonizes the walk time, and independent
Poisson coordinate counts give independent coordinate parities with
success probability $p(u)$. Since an odd-distance vertex has zero
probability at even walk times, no extra parity factor occurs.
The weight $w_\beta$ is integrable near $1/2$ and comparable to
$p^{-1-\beta}$ on $[0,b]$ for every fixed $b<1/2$.

Fix $\rho<b<1/2$ and take $j\le\rho D$. The integral on $[0,b]$
is comparable to the complete beta integral
$B(j-\beta,D-j+1)$. For the lower comparison, a beta random variable
with these parameters has mean at most $\rho$, so Markov's inequality
puts at least a fraction $1-\rho/b$ of its mass in $[0,b]$.
The identity
\[
 \binom Dj B(j-\beta,D-j+1)
 =\frac{\Gamma(D+1)\Gamma(j-\beta)}
        {\Gamma(j+1)\Gamma(D+1-\beta)}
 \asymp D^\beta j^{-1-\beta}
\]
gives the desired mass. On $[b,1/2]$, the binomial factor decreases
with $p$ and, after multiplication by $\binom Dj$, is at most
$\exp[-D\operatorname{KL}(\rho\Vert b)]$. The remaining integrable
weight makes this contribution exponentially small. This proves
\eqref{eq:level} and its stated extended range.

Sum \eqref{eq:poisson} over the odd distances above $r$. Dropping the
parity restriction and using the binomial Markov bound gives
\[
 \tau_r\le C D^{-\beta}\int_0^\infty
 u^{-1-\beta}\min\{1,Du/(r+1)\}\,du
 \le C(r+1)^{-\beta}.
\]
For the lower bound, sum \eqref{eq:level} over the odd $j$ in
$(r,2r+2]$, using a fixed cutoff strictly below $D/2$ as needed at
the endpoint. This yields $c r^{-\beta}$; bounded $r$ are absorbed
by a change of constants. Equation~\eqref{eq:ratio} follows from
\eqref{eq:level}, \eqref{eq:trace}, and \eqref{eq:paired}.

For the second moment, use paired multiplicities to write
$s_{2,r}=\sum_{j>r,\ j\text{ odd}}\binom Dj a_j^2$.
Up to an odd cutoff $b_D$ near $D/4$, its terms are at most
$Cj^{-2\beta-2}/\binom Dj$. Consecutive odd inverse binomial
coefficients have ratio
\[
 \frac{(j+1)(j+2)}{(D-j)(D-j-1)},
\]
uniformly less than a fixed constant below one in this range.
The sum is therefore bounded by a constant times its first term,
at $j=r+2$. Beyond $b_D$, monotonicity gives the bound
$a_{b_D}\tau_{b_D}\le C D^{-2\beta-1}/\binom D{b_D}$,
which is exponentially smaller than the claimed bound uniformly
for $r\le D/8$. This proves \eqref{eq:secondmoment}.
\end{proof}

Finally, two-step binomial ratios give, uniformly for odd $r\le D/8$,
\begin{equation}\label{eq:dimensionratios}
 \binom Dr\le D_r\le C\binom Dr,
 \qquad \frac{D_r}{\binom D{r+2}}\le C(r/D)^2.
\end{equation}
All constants used here can be fixed independently of a subsequently
chosen small positive limit $r/d\to\alpha$.

\begin{corollary}[Lower bounds used by the interpolation]\label{cor:lower-spectrum}
For sufficiently large $d$,
\[
 \lambda_0\ge c/d,\qquad
 \lambda_j\ge \frac{c d^{-1-\beta}}{\binom dj}\quad(0\le j\le .49d),
 \qquad v:=\lambda_{\min}(K)\ge c(d\log d)^{-\beta}.
\]
Also $v\to0$.
\end{corollary}
\begin{proof}
The first bound follows from the level estimate at full-cube degree one.
For odd reduced degree $j$, use $\lambda_j=a_j$ and
$\binom Dj/\binom dj=D/(D-j)$. For even $j$, use
$\lambda_j=a_{j+1}$ and
$\binom D{j+1}/\binom dj=D/(j+1)$.
The extended level bound with any fixed cutoff between $.49$ and $1/2$
then gives the stated uniform lower bound, including $j=0$.

For any odd full-cube degree $j\le D$, restrict the integral
\eqref{eq:poisson} to $u\ge\tfrac12\log(2D)$. In this range
$p(u)\ge(1-1/(2D))/2$ and $1-p(u)\ge1/2$, so
$p(u)^j(1-p(u))^{D-j}\ge c2^{-D}$. Integrating $u^{-1-\beta}$ gives
$a_j\ge c2^{-D}D^{-\beta}(\log D)^{-\beta}$.
Every reduced eigenvalue equals an odd $a_j$, and the matrix eigenvalues
are $M$ times the normalized ones. Thus $v\ge c(d\log d)^{-\beta}$.
Finally, for any fixed small positive cutoff fraction, the minimum
normalized eigenvalue is at most the average in its tail, whence
$v\le M\tau_r/(M-D_r)=O(d^{-\beta})$. This proves $v\to0$.
\end{proof}

\section{Random tail covariance matrices}\label{app:gram}
Let $Q$ be the covariance kernel retaining only degrees greater than an
odd cutoff $r$, and write $\tau=\tau_r$, $\ell=\lambda_{r+1}$, and
$s_2=s_{2,r}$. This section proves the sample event needed later, without
assuming that a Gaussian approximation remains valid after conditioning.
The entire argument concerns the exact limiting Gaussian covariance.

\begin{lemma}[Far-pair decay]\label{lem:far}
There are constants $c_0,C>0$ such that, when $r\le D/8$ is odd,
\[
 \max_{d/4\le\dist(x,y)\le3d/4}|Q(x,y)|
 \le\rho_r:=C\sqrt D\,e^{-r/8}+Ce^{-c_0D}.
\]
\end{lemma}
\begin{proof}
The full-cube Krawtchouk polynomial satisfies
\[
 \sum_{j=0}^D\mathsf K_j(h;D)z^j=(1+z)^{D-h}(1-z)^h.
\]
If $\delta D\le h\le(1-\delta)D$, put $b=\min(h,D-h)$.
On $|z|=t$ the generating function is bounded in absolute value by
\[
 (1+t^2)^b(1+t)^{D-2b}
 \le(1+t)^D\exp\!\left[-\frac{2\delta Dt}{(1+t)^2}\right].
\]
Cauchy's coefficient bound at $t=j/(D-j)$ and the uniform Stirling
lower bound for $\binom Dj$ therefore give, for $1\le j<D$,
\begin{equation}\label{eq:kraw}
 \frac{|\mathsf K_j(h;D)|}{\binom Dj}
 \le C\sqrt D\exp[-2\delta j(1-j/D)].
\end{equation}
For reduced-cube points at distance in $[d/4,3d/4]$, their augmented
full-cube distance lies in $[D/8,7D/8]$ for large $d$.
Because $r$ is odd, the exact tail is
$Q(x,y)=\sum_{j>r,\ j\text{ odd}}a_j\mathsf K_j(h;D)$.
For $r<j\le D-r$, \eqref{eq:kraw} gives a bound
$C\sqrt D e^{-r/8}$ after summation, since the total spectral mass
is one. For $j>D-r$, monotonicity and an odd cutoff $b_D$ near $D/4$
give
\[
 \sum_{j>D-r}\binom Dj a_j
 \le a_{b_D}\sum_{k<r}\binom Dk
 \le C D^{-\beta-1/2}
       \exp[-D(h(1/4)-h(1/8))+O(1)].
\]
The positive entropy gap absorbs the polynomial factor and proves the
claim. Adjust constants for bounded $d$.
\end{proof}

\begin{lemma}[A tail Gram deviation bound]\label{lem:gram-general}
Let $X_1,\ldots,X_n$ be independent uniform cube points, with
$1\ll n\le M/4$. Let
\[
 B_{\mathrm{near}}=\#\{y:\dist(x,y)\notin[d/4,3d/4]\}
 \le 2e^{d h(1/4)},\qquad q=2\log M.
\]
The count is independent of $x$. For every $z>0$,
\begin{align}
 &\mathbb P\{\|Q_{SS}-\tau I_n\|_{\mathrm{op}}>z\tau\}
 \notag\\
 &\quad\le\frac C{z^2}
 \left[\frac{q\rho_r}{\tau}
       +\frac{\sqrt{nq s_2}}{\tau}
       +\frac{n\ell}{\tau}
       +\frac{n(1+B_{\mathrm{near}})}M\right]^2
       +\frac{n^2(1+B_{\mathrm{near}})}M.
 \label{eq:gram-bound}
\end{align}
Here the sampled matrix includes repeated training locations.
\end{lemma}
\begin{proof}
Let $J=Q-\tau I_M$, a zero-diagonal Hermitian matrix. Form a deterministic
mask $J_{\mathrm{far}}$ retaining only pairs at distances in
$[d/4,3d/4]$. Since $Q$ is positive semidefinite with constant diagonal
$\tau$, every entry has absolute value at most $\tau$. Thus
\begin{align*}
 \|J_{\mathrm{far}}\|_{\max}&\le\rho_r,\\
 \|J_{\mathrm{far}}\|_{1\to2}^2&\le M s_2,\\
 \|J_{\mathrm{far}}\|_{\mathrm{op}}
 &\le M\ell+\tau+\tau B_{\mathrm{near}}.
\end{align*}
For the second inequality, Parseval gives squared column norm
$M s_2-\tau^2$ for $J$, and entrywise masking only decreases it.
For the third, $\|J\|\le M\ell+\tau$, while the removed matrix has
maximum absolute row sum at most $\tau B_{\mathrm{near}}$.
The entropy bound on $B_{\mathrm{near}}$ is the elementary binomial
tail bound on distances to a fixed cube point.

Let $R$ retain each of the $M$ coordinates independently with
probability $p=2n/M$. The zero-diagonal case of Tropp's random
principal-submatrix theorem \cite[Theorem 1.1]{tropp2008}, at moment
order two and $q=2\log M$, states
\[
 (\mathbb E\|RJR\|^2)^{1/2}
 \le C\left[q(\mathbb E\|RJR\|_{\max}^2)^{1/2}
       +\sqrt{pq}(\mathbb E\|JR\|_{1\to2}^2)^{1/2}
       +p\|J\|\right].
\]
Apply it to $J_{\mathrm{far}}$ and substitute the preceding bounds:
\begin{equation}\label{eq:bernoulli-gram}
 (\mathbb E\|R J_{\mathrm{far}}R\|^2)^{1/2}
 \le C[q\rho_r+\sqrt{nq s_2}+n\ell
               +(n/M)\tau(1+B_{\mathrm{near}})].
\end{equation}
The logarithm here is that of the original dimension $M$, not the
sample size.

To pass to a uniform subset of size $n$, use a uniform permutation of
all cube points, independently of $B\sim\operatorname{Bin}(M,2n/M)$.
Its first $B$ elements form the Bernoulli subset. On $B\ge n$, the
first $n$ principal submatrix is a compression of the larger one, so
its operator norm cannot be larger. The size-$n$ event is independent
of $B$, and $\mathbb P(B\ge n)\ge1/2$ for large $n$. Thus Markov's
inequality applied to \eqref{eq:bernoulli-gram} gives the same fixed-size
probability bound with a constant factor loss.

Conditional on no collisions, an independent size-$n$ sample is a
uniform size-$n$ subset, up to an irrelevant ordering. Collisions have
probability at most $n(n-1)/(2M)$. The probability of any sampled pair
violating the distance condition is at most $n^2B_{\mathrm{near}}/M$.
On their complement, the sampled matrices $J$ and $J_{\mathrm{far}}$
agree. Adding these exceptional probabilities proves
\eqref{eq:gram-bound}.
\end{proof}



\begin{thebibliography}{9}
\bibitem{flajolet2009} P. Flajolet and R. Sedgewick. \emph{Analytic Combinatorics}. Cambridge University Press, 2009. Chapter VI.
\bibitem{tropp2008} J. A. Tropp. Norms of random submatrices and sparse approximation. \emph{Comptes Rendus Math\'ematique} 346:1271--1274, 2008. \href{https://tropp.caltech.edu/papers/Tro08-Norms-Random-v3.pdf}{Author manuscript}, Theorem 1.1.
\end{thebibliography}
\end{document}
